\section*{Chapitre \ref{ch:iv1:estimation} --- Estimation}
\markboth{\protect\omTocMark{}{Réponses~: Estimation}}{\protect\omTocMark{}{Réponses~: Estimation}}

\begin{solution}{iq:iv1:estimation:1}
Une année compte environ $3.16 \times 10^7$ secondes, ce qui est proche de $\pi \times 10^7$. Les heures de négociation en comptent environ $5.9$ millions.

L'année compte $365.25 \cdot 86\,400 = 31\,557\,600$ secondes. Le raccourci $\pi \times 10^7$ est à un demi-pour-cent près, et il vaut la peine d'être retenu. Une année de négociation représente $252 \cdot 6.5 \cdot 3600 = 5\,896\,800$ secondes, si bien que le marché est ouvert environ $18.7~\%$ du temps, soit un peu moins d'un cinquième.

\iqfollow{Combien de minutes de négociation y a-t-il en une année~?}{$252 \cdot 390 = 98\,280$, soit environ $100\,000$.}
\iqwhy{Annualiser des chiffres intrajournaliers exige l'horloge de négociation. Un taux d'un événement par seconde de négociation représente environ six millions d'événements par an, et non trente millions.}
\end{solution}

\begin{solution}{iq:iv1:estimation:2}
Environ trois mois, avec un intervalle allant de six semaines à six mois.

La piscine contient $50 \cdot 25 \cdot 2 = 2500$ mètres cubes. Un tuyau d'arrosage remplit un seau de 10 litres en environ 30 secondes, soit 20 litres par minute, ou $1.2$ mètre cube par heure. La piscine demande donc $2500/1.2 \approx 2083$ heures, soit 87 jours sans interruption.

Le volume est exact, si bien que toute l'incertitude porte sur le débit. Entre 10 et 40 litres par minute, soit un facteur deux dans chaque sens, la durée passe de 174 jours à 43.

\iqwrong{«~Il faudrait quelques jours.~» Personne n'a d'intuition pour 2\,500 mètres cubes. C'est pourquoi on la convertit en litres, et le tuyau en un seau dont on se représente aisément le remplissage.}
\end{solution}

\begin{solution}{iq:iv1:estimation:3}
Un compteur sur 32 bits dure environ 72 minutes, un compteur sur 64 bits environ 585\,000 ans, et une remise à zéro quotidienne exige 35 bits.

Le premier vaut $2^{32}/10^6 \approx 4295$ secondes, en utilisant $2^{10} \approx 10^3$ pour obtenir d'un coup $4 \times 10^9$. Le second vaut $2^{64}/10^6 \approx 1.8 \times 10^{13}$ secondes, et une année compte environ $\pi \times 10^7$ secondes. Une journée compte $6.5 \cdot 3600 \cdot 10^6 = 2.34 \times 10^{10}$ ordres, valeur comprise entre $2^{34} \approx 1.7 \times 10^{10}$ et $2^{35} \approx 3.4 \times 10^{10}$. Donc 35 bits suffisent.

\iqwhy{Trente-deux bits paraissent largement suffisants jusqu'à un jour chargé. Un compteur qui déborde au bout d'une heure et quart tombe en panne en production, et non dans un test sur un petit échantillon. Choisir la largeur d'après ce calcul, et prévoir une marge pour que le débit puisse croître.}
\end{solution}

\begin{solution}{iq:iv1:estimation:4}
Pas si les jours sont indépendants~: 5~\% par jour passent à environ 79~\% par an, et 20~\% par an ne correspondent qu'à environ $1.3~\%$ par jour.

Avec des jours indépendants, les variances s'additionnent, si bien que la volatilité croît comme la racine carrée du temps\index{règle de la racine carrée du temps}. Le raccourci vaut $\sqrt{252} \approx 16$. Soit $5\% \cdot 16 \approx 80\%$ par an, soit $20\%/16 \approx 1.25\%$ par jour.

Pour que les deux affirmations soient vraies, les variations quotidiennes devraient presque s'annuler complètement. La variance annuelle effective ne serait alors que $(20/79.4)^2 \approx 6.3~\%$ de celle que donneraient des jours indépendants, ce qui supposerait un retour à la moyenne extrême.

\iqfollow{Quelques jours violents suffiraient-ils à les réconcilier~?}{Ils peuvent rendre la \emph{moyenne} trompeuse. Avec 248 jours à 1~\% et 4 jours de publication de résultats à 5~\%, la volatilité annuelle vaut $\sqrt{248 + 4 \cdot 25} \approx 18.7~\%$, de sorte que des mouvements de 5~\% sont possibles. Mais alors 5~\% n'est pas une journée ordinaire.}
\end{solution}

\begin{solution}{iq:iv1:estimation:5}
La fibre met 6 millisecondes dans un sens, une liaison radio en ligne droite en met 4, et un signal dans la fibre parcourt 20 centimètres par nanoseconde.

Dans la fibre, le signal se déplace à $200\,000$ km par seconde, donc $1200/200\,000$ secondes valent 6 ms. La radio dans l'air met $1200/300\,000$ secondes, soit 4 ms, ce qui fait gagner jusqu'à 2 ms dans chaque sens. Dans l'air, le signal parcourt environ 30 cm en une nanoseconde.

Un tracé réel de fibre est plus long que la ligne droite, si bien que le gain véritable est plus grand. Une liaison radio est moins fiable et transporte moins de données, ce qui explique pourquoi l'on utilise les deux.

\iqwhy{C'est le plancher de toute transaction entre deux villes~: aucun logiciel ni aucun matériel ne le bat. À l'intérieur d'un bâtiment, 20 cm par nanoseconde signifie que quelques mètres de câble supplémentaires coûtent 10 à 20 nanosecondes, ce qui peut annuler un gain logiciel durement acquis.}
\end{solution}

\begin{solution}{iq:iv1:estimation:6}
L'électricité coûte environ 657~\$ par an, soit environ un cinquième du coût annuel du serveur.

Une année compte environ 8\,760 heures, donc $0.5$ kW consomme $4380$ kWh, ce qui coûte 438~\$ à 10 cents le kWh. Le surcoût du bâtiment porte cela à $1.5 \cdot 438 = \$657$. La machine coûte $10\,000/4 = \$2500$ par an, si bien que l'électricité représente $657/(2500 + 657) \approx 21~\%$ du total.

\iqwrong{Oublier le surcoût, ou diviser par le seul coût de la machine. Une part est un rapport entre une partie et le tout, si bien que l'électricité doit figurer aux deux endroits.}
\iqwhy{Lorsqu'on propose d'économiser 10~\% de l'électricité d'un serveur, cela permet de savoir que l'économie représente environ 2~\% de la facture, avant tout changement du taux de panne ou de l'espace en baie.}
\end{solution}

\begin{solution}{iq:iv1:estimation:7}
Le cube a un côté d'environ 22 mètres, et il vaut environ 19~billions de dollars.

La densité vaut $19\,300$ kg par mètre cube, donc $2 \times 10^8$ kg occupent environ $10\,360$ mètres cubes. Comme $22^3 = 10\,648$, le côté vaut environ $21.8$ mètres. En onces, $2 \times 10^{11}/31.1 \approx 6.43 \times 10^9$, ce qui, à 3\,000~\$ l'once, donne environ $19.3$~billions de dollars.

\iqfollow{Par personne, avec 8 milliards d'habitants~?}{Environ 25 grammes chacun, soit moins d'une once.}
\iqwhy{La taille est la surprise~: toute cette valeur tiendrait dans un bâtiment modeste. C'est une bonne habitude pour les ordres de grandeur, car un chiffre qui semble énorme peut paraître petit une fois exprimé dans des unités que l'on peut se figurer.}
\end{solution}

\begin{solution}{iq:iv1:estimation:8}
Environ 2\,400, avec un intervalle allant de 1\,000 à environ 7\,000.

Raisonner à partir de la demande\index{estimation par la demande}. Un million de personnes, chacune avec environ 6 coupes par an, font 6 millions de coupes. Un coiffeur réalise environ une coupe toutes les 45 minutes, soit une dizaine de coupes par jour, sur environ 250 jours~: 2\,500 coupes par an. La division donne 2\,400.

Pour l'intervalle, on pousse conjointement les deux données dans le même sens. Avec 4 coupes par an et 4\,000 coupes par coiffeur, on obtient 1\,000. Avec 10 coupes par an et 1\,500 coupes par coiffeur, on obtient environ 6\,700. Leur moyenne géométrique\index{moyenne géométrique des bornes} est d'environ 2\,600, proche de la valeur centrale, ce qui est bon signe~: les deux bornes sont équilibrées.

La vérification s'appuie sur l'expérience d'une rue, et non sur celle d'une coupe de cheveux. La valeur centrale signifie un coiffeur pour 400 personnes. Un quartier de 2\,000 personnes compterait alors cinq coiffeurs, un ou deux salons de coiffure, ce qui paraît plausible.

\iqwrong{Donner un seul chiffre. Le recruteur veut voir quelles données restent incertaines, et l'on sait évaluer la largeur de la réponse.}
\end{solution}

\begin{solution}{iq:iv1:estimation:9}
Environ 650\,000, et les deux démarches donnent environ 690\,000 et 600\,000. Un intervalle de 500\,000 à 800\,000 est raisonnable.

Les sièges et les aménagements prennent environ un cinquième des $12 \cdot 2.5 \cdot 2 = 60$ mètres cubes, ce qui laisse 48.

La première démarche part du volume. Une balle de rayon $2.15$ cm a un volume $\tfrac43 \pi r^3 \approx 41.6$ centimètres cubes. Des sphères versées au hasard remplissent environ 60~\% de l'espace, si bien que le nombre vaut $48 \cdot 0.6/(41.6 \times 10^{-6}) \approx 692\,000$.

La seconde démarche les empile en grille cubique~: chaque balle occupe un cube de côté $4.3$ cm, soit $79.5$ centimètres cubes, et $48/(79.5 \times 10^{-6}) \approx 604\,000$. Une grille ne remplit que $\pi/6 \approx 52~\%$ de l'espace, ce qui explique ce résultat plus bas.

\iqwrong{Diviser le volume du bus par le volume d'une balle en négligeant les interstices. Cela surestime le résultat de plus de moitié. Même l'empilement le plus dense possible n'atteint que 74~\%, ce qui plafonne la réponse en dessous du million.}
\end{solution}

\begin{solution}{iq:iv1:estimation:10}
Environ un million et quart, avec un intervalle allant d'environ un demi-million à $2.5$ millions. Les deux démarches diffèrent d'environ 20~\%.

La première démarche part des passagers. Avec 8 milliards de personnes prenant chacune en moyenne un demi-vol par an, il y a 4 milliards de vols passagers. À environ $2.5$ heures chacun, cela représente $10^{10}$ heures-passager sur les $8760$ heures d'une année, si bien qu'en moyenne $1.14$ million de personnes sont en vol à un instant donné.

La seconde démarche part des avions. Avec environ 25\,000 avions de ligne, chacun volant 10 heures sur 24, environ $10\,400$ sont en vol à tout moment. À environ 130 passagers chacun, cela fait $1.35$ million.

C'est l'accord entre les deux démarches qui rassure, car elles ne partagent aucune donnée d'entrée. Le résultat signifie qu'environ une personne sur 7\,000 est en vol en ce moment. Les moyennes masquent un cycle journalier, si bien que l'instant considéré peut se situer un peu au-dessus ou au-dessous de la moyenne.

\iqwhy{Deux démarches indépendantes révèlent une donnée erronée. Si elles avaient différé d'un facteur dix, on saurait qu'une des deux est fausse, et cela est plus utile qu'un chiffre unique affiché avec assurance.}
\end{solution}

\begin{solution}{iq:iv1:estimation:11}
Environ 100\,000, avec un intervalle allant de 50\,000 à 200\,000.

Partir de la demande. Sur 300 millions d'habitants, on retient 200 millions de voitures. Si chacune fait le plein environ tous les 10 jours, cela représente 20 millions de pleins par jour. Une station de 8 pompes voit environ 200 pleins par jour, soit un toutes les cinq minutes sur une longue journée, si bien que le pays a besoin de $20\,000\,000/200 = 100\,000$ stations. Entre 400 et 100 pleins par jour, le nombre est compris entre 50\,000 et 200\,000.

Il faut maintenant le vérifier du côté de l'offre. Ce nombre signifie une station pour 3\,000 habitants. Une station gagne sur chaque plein~: à 40 litres et 5 cents de marge par litre, chaque plein rapporte 2~\$, soit 400~\$ par jour, ou environ 146\,000~\$ par an. Cela finance quelques salariés sur place, et la boutique ajoute davantage. Une station à 200 pleins par jour tout juste survit donc, ce qui suggère que le nombre ne peut guère être plus élevé.

\iqwhy{L'offre et la demande fournissent des vérifications de nature différente. Le côté demande indique ce dont on a besoin, le côté offre indique ce qui peut rester viable, et une véritable entreprise doit passer les deux épreuves.}
\end{solution}

\begin{solution}{iq:iv1:estimation:12}
D'un quart de l'estimation à quatre fois celle-ci, lorsque les facteurs sont indépendants et que les erreurs sont normales en logarithmes.

Passer aux logarithmes\index{erreurs lognormales}. Un produit de facteurs indépendants devient une somme d'erreurs indépendantes, et dans une somme les variances s'additionnent, si bien que les largeurs croissent comme la racine carrée du nombre de termes. Un intervalle à 90~\% d'un facteur 2 vaut $\pm \ln 2$ en logarithmes, donc $\sigma = \ln 2/1.645 \approx 0.42$. Avec quatre facteurs, la largeur double, pour atteindre $\pm 2 \ln 2 = \pm \ln 4$. Avec neuf, elle triplerait, pour atteindre un facteur 8.

Seize n'est atteint que si les quatre erreurs vont toutes dans le même sens à la fois, ce qui arrive si les facteurs sont parfaitement corrélés, puisqu'une seule erreur les entraîne toutes. Pour des facteurs indépendants, l'intervalle allant d'un seizième à seize contiendrait environ 99.9~\% des issues.

\iqfollow{Où se situe la moyenne du produit, par rapport à l'estimation centrale~?}{Au-dessus. La moyenne d'une loi lognormale vaut sa médiane multipliée par $e^{\sigma^2/2}$, avec $\sigma = 0.84$ ici, soit environ $1.43$ fois. Un chiffre central de Fermi est une médiane, et il sous-estime une moyenne.}
\iqwhy{Décomposer une quantité en de nombreux morceaux est ce qui rend les estimations de Fermi efficaces~: les erreurs indépendantes se compensent en partie. Mais cela ne fonctionne que si les morceaux ne partagent pas une même hypothèse.}
\end{solution}

\begin{solution}{iq:iv1:estimation:13}
Environ 1\,000 transactions par seconde. En pointe, le flux compte environ 610\,000 messages par seconde, soit environ 200 mégabits par seconde.

La valeur quotidienne échangée vaut $60 \times 10^{12} \cdot 2/252 \approx 476$~milliards de dollars. En divisant par 20\,000~\$, on obtient environ $23.8$ millions de transactions par jour. Une journée de négociation compte $6.5 \cdot 3600 = 23\,400$ secondes, soit environ 1\,020 transactions par seconde.

Avec 30 messages par transaction, la moyenne est d'environ $30\,500$ messages par seconde, et 20 fois ce chiffre donne $610\,000$. À 40 octets chacun, cela fait $24.4$ millions d'octets par seconde, soit $195$ mégabits par seconde.

\iqwrong{Dimensionner un système d'après la moyenne. Les taux de messages se concentrent fortement autour de l'ouverture, de la clôture et des annonces, et une conception qui tient la moyenne prend du retard précisément quand les données comptent le plus.}
\iqwhy{Un chiffre de ce type indique tout de suite quel matériel il faut~: une liaison de 1 gigabit a de la marge ici, mais seulement tant qu'on n'ajoute pas davantage de plateformes de négociation.}
\end{solution}

\begin{solution}{iq:iv1:estimation:14}
Environ 126~millions de dollars par an, avec un intervalle de 63 à 252 millions de dollars. Le chiffre central signifie que le teneur de marché en garde environ un cinquième de la demi-fourchette.

Cinq pour cent de 10 milliards font 500 millions d'actions par jour. À $0.1$ cent, qui est le milieu géométrique de $0.05$ et $0.2$, chaque jour rapporte 500\,000~\$, et sur 252 jours cela fait 126~millions de dollars. Les bornes varient directement avec les cents conservés.

La demi-fourchette vaut $0.5$ cent, ce que gagne un teneur de marché s'il exécute à son cours affiché et que le prix reste ensuite stable. En conserver $0.1$ cent signifie qu'environ $0.4$ cent est cédé à la sélection adverse\index{sélection adverse}, aux frais et au coût de détention de la position.

\iqwhy{Les revenus du market making semblent énormes exprimés en actions et minces par action. La fourchette n'est pas le profit~: l'essentiel sert à payer le fait d'être du mauvais côté face aux traders informés.}
\end{solution}

\begin{solution}{iq:iv1:estimation:15}
Environ $11.8$~To par an sans compression, soit environ $2.4$~To avec compression. Une journée sans compression prend environ 47 secondes à lire.

Une journée compte $5000 \cdot 10 \cdot 6.5 \cdot 3600 = 1.17 \times 10^9$ mises à jour. À 40 octets chacune, cela fait $46.8$ Go. Sur 252 jours, on arrive à $11.8$ To, et en divisant par 5 on obtient $2.36$ To. Au rythme d'un Go par seconde, $46.8$ Go prennent environ 47 secondes.

La moyenne masque un skew~: quelques centaines de titres actifs produisent l'essentiel des mises à jour, si bien qu'un titre très actif détient plusieurs fois sa part de $1/5000$ de la journée, soit environ 9 millisecondes de lecture, et un titre calme bien moins. Dans tous les cas, lire un seul titre est beaucoup plus rapide que parcourir toute la journée, si la disposition du fichier le permet.

\iqwhy{Le goulot d'étranglement de la recherche est rarement l'espace disque nécessaire pour une année. C'est le temps de parcours de nombreuses journées, et c'est pourquoi les bases de ticks partitionnent par date et par symbole.}
\end{solution}

\begin{solution}{iq:iv1:estimation:16}
Environ $2.7$ heures sur un cœur et environ 10 minutes sur 16 cœurs, moment où la lecture sur disque, de 6.5 minutes, devient presque le facteur limitant.

Les barres sont au nombre de $5000 \cdot 252 \cdot 20 \cdot 390 \approx 9.83 \times 10^9$. À 40 octets chacune, cela fait 393 Go, si bien que la lecture prend 393 secondes, soit environ $6.55$ minutes, quel que soit le nombre de cœurs. Le traitement à 1 microseconde par barre prend 9\,828 secondes, soit $2.73$ heures, sur un cœur.

Les deux se recouvrent, si bien que sur un cœur le traitement est de loin la contrainte. Sur 16 cœurs, le traitement prend 614 secondes, toujours plus lent que la lecture sur disque. Les deux durées deviennent égales pour $9828/393 = 25$ cœurs, et au-delà, ajouter des cœurs ne sert à rien. Il faut un disque plus rapide, ou lire moins de données.

\iqwhy{Le résultat utile est le rapport, non les $2.7$ heures. Il faut savoir quelle ressource est limitante avant de dépenser de l'argent~: à 25 cœurs, en ajouter n'apporte plus rien, et stocker les données par colonne, de sorte qu'une passe n'en lise qu'un dixième, profite à tout le monde.}
\end{solution}

\begin{solution}{iq:iv1:estimation:17}
Environ $12.6$~milliards de dollars par an, soit environ $2.1$~points de base de la valeur du marché.

Chaque transaction a un côté qui traverse la fourchette, et ce côté paie la moitié de la fourchette affichée par rapport au prix milieu, soit ici 1 point de base. Le péage quotidien vaut donc $500 \times 10^9 \cdot 10^{-4} = 50$~millions de dollars, et sur 252 jours cela fait \$$12.6$~milliards de dollars. Rapporté aux 60~billions de dollars que représente la valeur totale du marché, cela fait $2.1$ points de base par an.

\iqwrong{Facturer la fourchette complète, ce qui donne $25.2$~milliards de dollars. Un preneur de liquidité paie la moitié de la fourchette par rapport au prix milieu, et payer la fourchette complète supposerait que la même partie achète puis vende. Une autre erreur consiste à facturer les deux côtés de chaque transaction~: le côté passif gagne ce que paie le côté qui traverse.}
\iqwhy{C'est un plafond utile pour ce que les activités de capture de la fourchette peuvent se partager, avant l'impact de marché et la sélection adverse. Cela montre aussi combien le péage est faible pour un détenteur de long terme.}
\end{solution}

\begin{solution}{iq:iv1:estimation:18}
Non. Cela implique 1~\% de la valeur échangée totale du titre chaque jour, ce qui dépasse de loin tout edge plausible.

Le chiffre décisif est le profit rapporté au volume~: $1$~million de dollars sur 100~millions de dollars représente 100 points de base sur chaque dollar échangé par quiconque ce jour-là, même si la stratégie a réalisé toutes les transactions. Une part réaliste de 10~\% signifie 10~millions de dollars échangés, si bien qu'elle rapporte 10~\% par dollar échangé. Cela représente cinq fois l'écart type quotidien du titre, de 2~\%, chaque jour.

Un seul jour de ce type est déjà une valeur aberrante, et y parvenir chaque jour sans aucune perte ne peut pas relever de l'habileté au sens ordinaire. Cela suggère une valorisation périmée, une erreur de données, un modèle d'exécution irréaliste ou une mauvaise attribution.

\iqfollow{Que demanderiez-vous ensuite~?}{Comment les exécutions ont été modélisées, si le P\&L était valorisé à des cours auxquels on pouvait réellement négocier, et quelle était la taille de la position à la clôture. Ces questions trouvent l'erreur plus vite que n'importe quel débat sur le signal.}
\iqwhy{Diviser le profit par le volume est le moyen le plus rapide de vérifier un historique de performance, et il n'exige aucun détail sur la stratégie.}
\end{solution}

\begin{solution}{iq:iv1:estimation:19}
En moyenne environ $1.5$ molécules, si bien que la probabilité d'en avoir au moins une est d'environ 77~\%.

Commençons par l'atmosphère. La pression est le poids de l'air au-dessus d'un mètre carré, si bien que chaque mètre carré porte $10^5/10 = 10^4$ kg. La surface de la Terre vaut $4\pi R^2 \approx 5.15 \times 10^{14}$ mètres carrés, si bien que l'atmosphère pèse environ $5.15 \times 10^{18}$ kg. À 29 grammes par mole, cela fait environ $1.77 \times 10^{20}$ moles, soit environ $1.06 \times 10^{44}$ molécules.

Une respiration représente $0.5/24 \approx 0.0208$ mole, soit $n \approx 1.25 \times 10^{22}$ molécules. Le souffle de César, réparti uniformément, représente une part $n/N$ de l'atmosphère, et chacune des $n$ molécules inhalées a cette probabilité d'être l'une des siennes, si bien que la moyenne vaut $n^2/N \approx 1.46$.

Le décompte est une somme de nombreux événements indépendants très rares, il suit donc une loi de Poisson\index{approximation de Poisson}, et $P(\text{au moins un}) = 1 - e^{-1.46} \approx 0.77$.

\iqwrong{Supposer que la réponse doit être minuscule ou gigantesque, puisqu'une respiration est si petite et l'air si vaste. Les deux très grands nombres se compensent presque~: $n$ intervient au carré, et $n^2$ se trouve être proche de $N$.}
\iqfollow{Quelle hypothèse est la plus fragile~?}{Un mélange uniforme, et le fait que les molécules restent inchangées. Sur 2\,000 ans, l'air se mélange bien, mais une grande partie du dioxyde de carbone et une partie de l'oxygène sont absorbées puis remplacées. L'azote est la partie la plus stable, si bien que l'estimation constitue une borne supérieure.}
\end{solution}

\begin{solution}{iq:iv1:estimation:20}
$52!$ vaut environ $8 \times 10^{67}$, soit environ $7 \times 10^{17}$ fois le nombre d'atomes de la Terre. La machine aurait vu $4.4 \times 10^{26}$ permutations, et une chance sur deux de trouver une coïncidence demande environ $10^{34}$ mélanges.

Utiliser la formule de Stirling\index{formule de Stirling}~: $\ln n! \approx n \ln n - n + \tfrac12 \ln(2\pi n)$. Ici $52 \ln 52 \approx 205.5$~; retrancher 52 et ajouter $2.9$ donne $156.4$. En divisant par $\ln 10 \approx 2.303$, on obtient $67.9$, donc $52! \approx 8 \times 10^{67}$. La forme approchée $n \ln n - n$ seule donne 66.6, ce qui montre pourquoi le dernier terme compte à cette précision.

La Terre renferme $6 \times 10^{27}$ grammes, soit $2 \times 10^{26}$ moles d'atomes, soit $1.2 \times 10^{50}$ atomes. Le rapport vaut environ $6.7 \times 10^{17}$. La machine a passé en revue $10^9 \cdot 4.4 \times 10^{17} = 4.4 \times 10^{26}$ permutations, ce qui laisse un facteur de $1.8 \times 10^{41}$.

Pour une coïncidence, utiliser l'estimation du problème des anniversaires\index{problème des anniversaires}. Avec $k$ mélanges, il y a environ $k^2/2$ paires, chacune coïncidant avec la probabilité $1/52!$, si bien qu'une chance sur deux demande $k \approx \sqrt{2 \ln 2 \cdot 52!} \approx 1.06 \times 10^{34}$.

\iqwhy{Cela explique pourquoi les identifiants aléatoires avec assez de bits sont sûrs en pratique~: la racine carrée du décompte des anniversaires fixe la taille, si bien que 128 bits sont confortables là où 64 ne le sont pas.}
\end{solution}

\begin{solution}{iq:iv1:estimation:21}
Cinq bonnes réponses ou moins ont une probabilité d'environ $0.16~\%$. Élargir les intervalles d'un facteur d'environ $2.44$ en logarithmes, si bien que «~à un facteur 3 près~» doit devenir «~à un facteur d'environ 15 près~».

Le score suit une loi binomiale de paramètres $n = 10$ et $p = 0.9$, et $P(X \le 5) \approx 0.00163$. On n'a donc pas simplement eu de malchance~: les intervalles sont trop étroits. Même avec si peu de données, un intervalle exact à 95~\% pour la couverture réelle va d'environ 19~\% à 81~\%, et n'atteint pas 90~\%.

Pour fixer les idées, supposons que les erreurs en logarithmes suivent une loi normale d'écart type vrai $\sigma$. Un intervalle réel à 90~\% vaut $\pm 1.645\sigma$. La couverture obtenue est d'environ 50~\%, ce qui correspond à $\pm 0.674\sigma$. Ce qui a été pris pour $1.645$ fois le $\sigma$ supposé vaut en réalité $0.674$ fois le vrai, et le rapport vaut $1.645/0.674 \approx 2.44$. Un facteur 3 vaut $\ln 3$ en logarithmes, et $e^{2.44 \ln 3} \approx 14.6$.

\iqwrong{Élargir selon le rapport des couvertures, $90/50$. La couverture n'est pas linéaire en largeur, car les queues d'une loi normale sont fines, si bien qu'il faut une bonne part de largeur supplémentaire pour attraper les 40~\% restants.}
\iqwhy{La surconfiance est le résultat habituel de ce type de question, et les recruteurs qui demandent un intervalle guettent précisément ce travers. Sur un desk, cette calibration fixe les limites de position~: se tromper en dehors de son intervalle la moitié du temps est un problème de risque.}
\end{solution}

\begin{solution}{iq:iv1:estimation:22}
Environ $4.1~\%$ si la loi est normale, et environ 22~\% sous la loi de Student $t$ à 3 degrés de liberté. Une baisse d'environ 20~\% rejette le modèle normal, et elle n'a rien de surprenant sous le modèle à queues épaisses.

Il y a $25\,200$ jours, donc il faut chercher le mouvement dépassé un de ces jours, ce qui correspond à une queue bilatérale de $1/25\,200$. Pour une loi normale, $P(|Z| > z) = 1/25\,200$ donne $z \approx 4.11$, soit $4.1~\%$.

Une loi de Student $t$ à 3 degrés de liberté a une variance égale à 3, il faut donc la multiplier par $1/\sqrt{3}$ pour obtenir une volatilité quotidienne de 1~\%. Sa queue ne décroît qu'en loi de puissance\index{queues épaisses}, et le même seuil de queue vaut environ 22~\%. La médiane du maximum sur 100 ans est un peu plus élevée, environ $4.2\%$ et environ 25~\%, puisqu'elle se situe à une queue de $\ln 2/25\,200$.

Une baisse de 20 écarts types sous une loi normale a une probabilité d'environ $3 \times 10^{-89}$ par jour, ce qui représente une attente d'environ $10^{86}$ années. Pourtant, le Dow Jones Industrial Average a perdu $22.6~\%$ en une seule séance en octobre 1987, donc c'est ce modèle qu'il faut abandonner.

\iqwhy{Les estimations fondées sur un modèle normal paraissent précises, et elles peuvent se tromper d'un facteur cinq dans la queue, là où le capital est nécessaire. Il faut demander l'hypothèse sur les queues dès qu'on avance un pire scénario.}
\end{solution}

\begin{solution}{iq:iv1:estimation:23}
Le livre peut atteindre $6.25$~milliards de dollars. Il rapporte environ 315 millions de dollars par an, soit environ 5~\% du livre.

La limite quotidienne de négociation vaut $1000 \cdot 50 \times 10^6 \cdot 5~\% = \$2.5$ milliards. Il faut maintenant compter les transactions avec soin. Pour garder un livre de taille $G$ avec une période de détention de 5 jours, chaque jour on clôture $G/5$ et on en ouvre un autre $G/5$, donc on négocie $2G/5$ par jour. Poser $2G/5 = 2.5$ milliards donne $G = 6.25$ milliards.

Le volume échangé vaut $2.5$ milliards par jour, soit 630 milliards de dollars par an. À 5 points de base sur chaque dollar négocié, cela fait 315 millions de dollars, soit $5.04~\%$ sur $6.25$ milliards de dollars.

\iqwrong{Multiplier la limite quotidienne par la période de détention pour obtenir $12.5$ milliards de dollars. Cela ne compte que la phase d'ouverture. Une position doit aussi être clôturée, et chacune de ces transactions consomme le même budget de volume.}
\iqfollow{Sans changer le signal, quel est le moyen le plus simple d'augmenter la capacité~?}{Conserver la position plus longtemps, si le signal persiste. Avec une période de détention de 10 jours, les mêmes transactions soutiennent un livre deux fois plus grand, même si le rendement par dollar du livre baisse, sauf si le signal croît dans les mêmes proportions.}
\iqwhy{La capacité se décide en général par ce calcul, et non par le modèle. Les investisseurs la demandent tôt, car une participation de 5~\% sur cinq jours fixe la taille maximale que le fonds pourra jamais atteindre.}
\end{solution}

\begin{solution}{iq:iv1:estimation:24}
Le temps d'attente moyen vaut $1.6$ microsecondes. La rafale se termine avec une file de 250 messages, qui se vide environ 333 microsecondes plus tard.

La capacité du gestionnaire vaut $1/0.8 = 1.25$ message par microseconde, donc sa charge vaut $\rho = 0.8$. Avec des arrivées aléatoires et un temps de service fixe, la file est une file M/D/1\index{file M/D/1}, et la formule de Pollaczek-Khinchine donne
\[
W = \frac{\rho\, s}{2(1 - \rho)} = \frac{0.8 \cdot 0.8}{2 \cdot 0.2} = 1.6.
\]
Le temps total dans le système vaut $2.4$ microsecondes. Si les temps de service étaient exponentiels, l'attente doublerait pour atteindre $3.2$. Le terme $1/(1 - \rho)$ est l'avertissement~: à une charge de 90~\%, l'attente est plus que doublée.

Pendant la rafale, l'excès vaut $1.5 - 1.25 = 0.25$ message par microseconde pendant 1\,000 microsecondes, soit 250 messages, ce qui représente un retard de 200 microsecondes. Ensuite, la file se réduit à $1.25 - 0.5 = 0.75$ par microseconde, donc elle se vide en $250/0.75 \approx 333$ microsecondes. Autrement dit, une rafale de 1 ms gêne pendant environ $1.3$ ms.

\iqwhy{La charge moyenne masque ce risque. Un gestionnaire à 40~\% de charge moyenne peut tout de même accuser un retard de 200 microsecondes lors d'une rafale, exactement comme le marché bouge. Il faut dimensionner pour la rafale, et non pour la moyenne.}
\end{solution}

\begin{solution}{iq:iv1:estimation:25}
Non~: les coûts seuls représentent environ 42 points de base par dollar négocié, contre environ 6 points de profit net annoncés. Avec un signal de 10 points de base, le meilleur livre vaut environ 490 millions de dollars et rapporte environ 83 millions de dollars par an.

Le fonds négocie $2 \cdot 20$ milliards de dollars par jour, soit 8~\% du marché, donc $q = 0.08$, et le coût par dollar négocié vaut $\sigma\sqrt{q} = 150 \cdot 0.283 \approx 42.4$ points de base. La prétention est de 30~\% de 20 milliards de dollars, soit 6 milliards de dollars par an. Réparti sur 252 jours et 40 milliards de dollars par jour, cela fait environ $5.95$ points de base nets. Le signal devrait donc rapporter environ $48$ points de base sur chaque dollar négocié, chaque jour, sur un volume d'environ un douzième du marché. Avec un signal réel de 10 points de base, cette taille perd environ 130 millions de dollars par jour.

Optimisons maintenant. Le profit par jour, avec $q$ la part du fonds et $V$ la valeur du marché, vaut
\[
\pi(q) = qV\bigl(a - \sigma\sqrt{q}\bigr),
\]
et annuler la dérivée donne $a = \tfrac32\sigma\sqrt{q}$, soit $q^* = (2a/3\sigma)^2 \approx 0.00198$. On négocie donc environ 988 millions de dollars par jour, sur un livre d'environ 494 millions de dollars. Un coût égal aux deux tiers du signal est la signature de la loi en racine carrée\index{impact en racine carrée}. Le net vaut $a/3 \approx 3.33$ points de base, soit environ 329\,000 dollars par jour et 83 millions de dollars par an, environ 17~\% du livre.

\iqwrong{S'arrêter quand les coûts égalent le signal. Le profit vaut alors zéro~: l'optimum se situe bien avant, là où le dernier dollar atteint tout juste l'équilibre mais où le dollar moyen rapporte encore.}
\iqwhy{C'est le contrôle à appliquer à toute prétention importante, et ce résultat explique pourquoi les bons signaux soutiennent de petits livres~: le rendement baisse quand le livre grandit, alors que le profit en dollars augmente beaucoup plus lentement.}
\end{solution}
